% ID: SEL-D0052-OA-E03
% ESTADO: cotejada. Segunda pasada visual y borrador compilado revisados.
% Fuente: Selectividad/sel_2013_matematicas_aplicadas/septiembre/reserva b/EXAMEN 4-ANDALUCÍA.pdf
% SHA-256: 6201aa5a68f0d1302cfd697751f74c4c73f655f0d9b1ae98ae8960207dbd5225
% Página: 1; opción A; ejercicio original 3.

Un Centro de Salud propone dos terapias, A y B, para dejar de fumar. De las personas que acuden al Centro para dejar de fumar, el 45\% elige la terapia A, y el resto la B. Después de un año el 70\% de los que siguieron la terapia A y el 80\% de los que siguieron la B no han vuelto a fumar.

Se elige al azar un usuario del Centro que siguió una de las dos terapias:
\begin{enumerate}[label=\alph*)]
\item \textbf{(1 punto)} Calcule la probabilidad de que después de un año no haya vuelto a fumar.
\item \textbf{(0.75 puntos)} Si transcurrido un año esa persona sigue sin fumar, calcule la probabilidad de que hubiera seguido la terapia A.
\item \textbf{(0.75 puntos)} Si transcurrido un año esa persona ha vuelto a fumar, calcule la probabilidad de que hubiera seguido la terapia A.
\end{enumerate}
